\section{Модель нарушителя и угроз}
\label{sec:adversary}

\subsection{Участники системы и их роли}
% ЧЕРНОВИК
Владельцы данных, три вычислительных сервера, владелец модели, получатель
результата, администратор.

\subsection{Классификация нарушителей}
Внешний и внутренний; получестный и активный; статический и адаптивный.
Порог компрометации для ABY3~--- один сервер из трёх~\cite{mohassel2018aby3,araki2016}.

\subsection{Возможности и цели нарушителя}

\subsection{Допущения безопасности}
Защищённые каналы между серверами, независимое администрирование серверов,
корректная генерация случайности.

\subsection{Формальная постановка}
Безопасность определяется в парадигме <<реальный мир~--- идеальный мир>>.

\subsection{Перечень актуальных угроз}
% Оформить по Методике оценки угроз ФСТЭК России (2021).

\subsection{Остаточные риски}
Методика моделирования атак на модели машинного обучения~\cite{podtopelny2024}.
